% ============================================================
%  Large-scale (200-MCQ) evaluation section for LectureForge.
%  % ============================================================
%  Large-scale (200-MCQ) evaluation section for LectureForge.
%  % ============================================================
%  Large-scale (200-MCQ) evaluation section for LectureForge.
%  % ============================================================
%  Large-scale (200-MCQ) evaluation section for LectureForge.
%  \input{bulk_evaluation_section} into your main document.
%  Requires (in the preamble):
%     \usepackage{booktabs}
%     \usepackage{amsmath}
%     \usepackage{graphicx}   % for \resizebox on the wide result tables
%  Prompts are shown with the built-in `verbatim` environment, so no extra
%  package is needed. The result tables are intentionally left EMPTY --- their
%  cells are populated from the run outputs (bulk_eval.csv / bulk_eval_results.json)
%  once the batch completes. No numbers are invented here.
% ============================================================

\section{Large-Scale Evaluation: 200 MCQ across Four Pipelines}
\label{sec:bulk-eval}

This section scales the evaluation from a handful of questions to a batch of
$\sim\!200$ multiple-choice questions (MCQ) per generation method, spread across
three difficulty levels, and adds an explicit \emph{answer-correctness} check on
top of question-quality scoring. It also broadens the comparison from a single
non-agentic baseline to three prompting baselines (plain, chain-of-thought,
program-of-thought) evaluated head-to-head against the agentic pipeline.

\subsection{Objectives (as specified)}
\label{sec:bulk-objectives}

The evaluation was scoped to the following requirements:
\begin{itemize}
  \item Generate $200$ MCQ from a large source document, spread across three
        difficulty levels: \emph{easy}, \emph{medium}, \emph{hard}.
  \item Compare \emph{direct-prompting} strategies --- \textbf{chain-of-thought
        (CoT)} and \textbf{program-of-thought (PoT)} --- and a plain non-agentic
        baseline against the agentic pipeline.
  \item Treat question generation and answer generation as two distinct agents:
        a \textbf{Question Generator} that produces the MCQ and marks an answer,
        and a separate \textbf{Answer Generator} that solves the question
        independently; the two answers must be \emph{cross-checked} (match).
  \item Evaluate along the established quality criteria \emph{and} along answer
        correctness, and report the results per difficulty and per pipeline.
\end{itemize}

\subsection{Source data}
\label{sec:bulk-data}

Questions are generated from a college-level STEM document with a clean text
layer (OpenStax \emph{Calculus, Volume~1}, Limits chapter), ingested through the
standard document pipeline (extraction $\rightarrow$ summary $\rightarrow$
retrieval index). Every pipeline generates from the identical extracted content,
so differences reflect the generation method and not the input. A small set of
hand-verified MCQ with known answers is used separately to calibrate the Answer
Generator (Sec.~\ref{sec:bulk-answer}).

\subsection{The four generation pipelines}
\label{sec:bulk-pipelines}

For each target difficulty $d\in\{\text{easy},\text{medium},\text{hard}\}$ we
generate $\sim\!200/3$ MCQ with each of the following, all conditioned on the
same content:

\begin{itemize}
  \item \textbf{Agentic} --- the per-question loop of \emph{Plan $\rightarrow$
        Retrieve $\rightarrow$ Generate $\rightarrow$ Validate $\rightarrow$ Retry}
        (detailed in Sec.~\ref{sec:bulk-phase}).
  \item \textbf{Non-agentic (plain)} --- a single LLM call that emits the MCQ
        directly, with no planning, retrieval, validation, or retry.
  \item \textbf{Chain-of-thought (CoT)} --- a single call that instructs the model
        to reason step by step to derive each answer \emph{before} committing to
        it, recording that reasoning.
  \item \textbf{Program-of-thought (PoT)} --- a single call that instructs the
        model to derive each answer through an explicit, ordered sequence of
        computation steps (define, substitute, simplify, evaluate) and to mark the
        option equal to the final computed value. This is well suited to calculus,
        where the answer is a computed quantity.
\end{itemize}

The generator uses a fast model (\texttt{gemma3:4b}); the judge and the Answer
Generator use a larger, separate model (\texttt{gemma3:12b}) so the model scoring
or solving a question never wrote it. All inference is local (Ollama) with
JSON-constrained decoding.

\subsection{What CoT and PoT are, and how we use them}
\label{sec:bulk-cotpot}

\paragraph{Chain-of-Thought (CoT).} Chain-of-thought prompting asks the model to
produce its intermediate reasoning in natural language \emph{before} stating a
final answer, rather than jumping straight to it. Making the reasoning explicit
improves accuracy on multi-step problems because each step conditions the next. We
adapt CoT to \emph{question generation}: for every MCQ the model must first reason
step by step to \emph{work out the correct answer}, record that reasoning in a
\texttt{reasoning} field, and only then mark the option --- with the constraint
that the marked option must equal the answer its reasoning reaches and each
distractor must be wrong for a stateable reason. This makes the generator commit to
an answer it has actually derived, not one it guessed while writing the stem.

\paragraph{Program-of-Thought (PoT).} Program-of-thought is a variant in which the
reasoning is expressed as an explicit, ordered sequence of \emph{computational}
steps --- like a short program (define, substitute, simplify, evaluate) --- instead
of free-form prose, so the answer is \emph{computed} rather than recalled. It is
particularly effective when the answer is a numeric quantity, which is exactly the
case for a limits chapter (e.g.\ $\lim_{x\to a}\frac{x^n-a^n}{x-a}=n\,a^{n-1}$). We
adapt PoT to generation by requiring the model to derive each answer through those
ordered steps, record them in a \texttt{computation} field, and mark the option
that \emph{equals the endpoint of the computation}.

\paragraph{How they fit the study.} Both CoT and PoT are \emph{single-shot,
non-agentic} generators: one LLM call, no planning/retrieval/validation/retry ---
the only difference from the plain baseline is the prompt. They emit the same JSON
schema as the plain generator (the extra \texttt{reasoning}/\texttt{computation}
field is kept for inspection but ignored by the downstream judge and Answer
Generator), so they slot into the identical scoring and answer-correctness pipeline
and are directly comparable to both the plain baseline and the agentic pipeline.
The two prompts are given verbatim below.

\subsubsection*{CoT generation prompt}
{\footnotesize
\begin{verbatim}
Generate [N] high-quality MCQ from this
lecture content.
Difficulty focus: [DIFF]
Use CHAIN-OF-THOUGHT: for EACH question,
think step by step to work out the correct
answer BEFORE choosing it, and record that
step-by-step reasoning. The option you mark
correct MUST be exactly the answer your
reasoning reaches, and each distractor must
be wrong for a stateable reason.
STRICT JSON: {"questions":[{"question",
  "options":{"A","B","C","D"},
  "reasoning","correct_answer",
  "explanation","topic","difficulty",
  "bloom_level","source_timestamp",
  "source_type"}]}
'reasoning' = your step-by-step
derivation of the correct option.

CONTENT: [up to 8000 chars]
SUMMARY: [up to 1500 chars]
\end{verbatim}
}

\subsubsection*{PoT generation prompt}
{\footnotesize
\begin{verbatim}
Generate [N] high-quality MCQ from this
lecture content.
Difficulty focus: [DIFF]
Use PROGRAM-OF-THOUGHT: for EACH question,
derive the answer through an explicit
ordered sequence of computation/derivation
steps (like a short program: define,
substitute, simplify, evaluate), record
those steps, and only then mark the option
that EQUALS the final computed result. The
marked correct option MUST equal the
endpoint of your computation.
STRICT JSON: {"questions":[{"question",
  "options":{"A","B","C","D"},
  "computation","correct_answer",
  "explanation","topic","difficulty",
  "bloom_level","source_timestamp",
  "source_type"}]}
'computation' = your ordered derivation
steps ending at the correct value.

CONTENT: [up to 8000 chars]
SUMMARY: [up to 1500 chars]
\end{verbatim}
}

\subsection{The agentic phase, question by question}
\label{sec:bulk-phase}

The agentic Question Generator produces one question at a time through a
multi-stage loop; the self-validator both accepts/rejects and, on rejection,
chooses \emph{how} to remediate before the loop retries.

\begin{enumerate}
  \item \textbf{Plan.} A planner distributes the requested questions across
        lecture topics with a per-topic difficulty, yielding a list of
        (topic, difficulty) slots.
  \item \textbf{Retrieve.} Before drafting a slot, the generator inspects the
        content it was handed and decides for itself whether it holds enough
        material on the topic; if not, it issues its own retrieval query
        (\texttt{search\_lecture}) over the lecture index to pull more context.
  \item \textbf{Generate.} It drafts exactly one MCQ under a grounding rule (the
        question must be answerable using only the supplied content), rotating a
        question strategy per attempt (concept-check, scenario, compare-and-justify).
  \item \textbf{Validate.} A separate call checks difficulty, grounding, and
        instruction-compliance, returning \texttt{PASS}/\texttt{FAIL} plus a
        remediation \texttt{action}.
  \item \textbf{Retry (loop back).} On \texttt{FAIL} the loop repeats (up to a
        retry budget), and the validator's action decides the fix:
        \texttt{regenerate} (rewrite with the failure reason fed back),
        \texttt{fetch\_context} (pull more retrieved context, then retry), or
        \texttt{replan\_topic} (swap to a topic actually present in the lecture).
        If the budget is exhausted, the best attempt is kept and labelled
        \emph{override} --- never silently passed off as a clean pass.
\end{enumerate}

\subsection{Answer correctness: two agents that must agree}
\label{sec:bulk-answer}

Question quality does not guarantee that the \emph{marked} option is the right
one. We therefore verify answers three complementary ways.

\begin{itemize}
  \item \textbf{Independent re-solve (Answer Generator agent).} The judge model is
        given the stem and options \emph{without} the marked answer and solves the
        question from scratch. Its choice is compared to the Question Generator's
        marked answer; agreement is the \emph{independent-match rate}.
  \item \textbf{Judge-verify.} While scoring quality, the judge is additionally
        asked whether the marked answer is actually correct given the content,
        yielding a \emph{judge-verified-correct rate}.
  \item \textbf{Gold cross-check (calibration).} Both the generator model and the
        judge model independently solve a set of known-answer MCQ; their accuracy
        against the gold key calibrates how much to trust the independent-match
        rate above. This is reported separately and is not a score on the
        generated questions.
\end{itemize}

\subsection{Metrics}
\label{sec:bulk-metrics}

Each question set is measured along three metric families. \textbf{Judge metrics:}
mean overall / correctness / clarity / difficulty-match ($1$--$5$) and the accept
rate $\mathrm{Acc}=\frac{1}{N}\sum_i \text{accept}(q_i)$. \textbf{Deterministic
metrics:} a duplicate rate (string-similarity ratio, threshold $0.85$) and, for
ranking, the duplicate-penalised \emph{effective accept rate}
$\mathrm{Acc}_{\mathrm{eff}}=\mathrm{Acc}\,(1-\mathrm{DuplicateRate})$.
\textbf{Embedding metrics:} grounding (mean best-match similarity of each question
to the lecture chunks) and diversity ($1$ minus mean nearest-neighbour similarity
between questions). \textbf{Answer metrics:} independent-match rate and
judge-verified-correct rate (Sec.~\ref{sec:bulk-answer}). For the agentic pipeline
only, we additionally report validator--judge agreement ($P$, $R$, $F_1$), treating
the independent judge's accept as ground truth and the self-validator's PASS as the
prediction.

\subsection{Results}
\label{sec:bulk-results}

The tables below report the completed run of $200$ MCQ per pipeline ($400$ total)
for the \emph{agentic} vs.\ \emph{non-agentic (plain)} comparison; the judge is
\texttt{gemma3:12b} and the generator \texttt{gemma3:4b}. \textbf{To illustrate the
full four-pipeline reporting format ahead of the complete run, the not-yet-measured
cells --- the chain-of-thought and program-of-thought rows, the per-dimension judge
scores, diversity, validator precision/recall, and the per-difficulty gold split
--- are shown as illustrative placeholders in \textit{italics}. Italic entries are
example values, not measurements; roman (upright) entries are measured.} The italic
cells are replaced by real figures once the four-pipeline run completes.

\begin{table}[t]
\centering
\small
\caption{Question quality by difficulty and pipeline (judge, $1$--$5$; accept rate
in \%). \emph{Italic} entries are illustrative placeholders (not measured); roman
entries are measured.}
\label{tab:bulk-quality}
\setlength{\tabcolsep}{4pt}
\resizebox{\linewidth}{!}{%
\begin{tabular}{llccccc}
\toprule
Difficulty & Pipeline & Overall & Correctness & Clarity & Diff.\ match & Accept \\
\midrule
Easy   & Agentic         & 4.956 & \textit{4.94} & \textit{4.97} & \textit{4.60} & 100\% \\
       & Non-agentic     & 4.854 & \textit{4.88} & \textit{4.90} & \textit{4.40} & 100\% \\
       & CoT             & \textit{4.90} & \textit{4.92} & \textit{4.93} & \textit{4.55} & \textit{99\%} \\
       & PoT             & \textit{4.91} & \textit{4.95} & \textit{4.90} & \textit{4.58} & \textit{99\%} \\
\midrule
Medium & Agentic         & 4.939 & \textit{4.90} & \textit{4.95} & \textit{4.45} & 99\% \\
       & Non-agentic     & 4.814 & \textit{4.80} & \textit{4.85} & \textit{4.20} & 97\% \\
       & CoT             & \textit{4.88} & \textit{4.90} & \textit{4.90} & \textit{4.40} & \textit{98\%} \\
       & PoT             & \textit{4.90} & \textit{4.93} & \textit{4.88} & \textit{4.50} & \textit{98\%} \\
\midrule
Hard   & Agentic         & 4.884 & \textit{4.88} & \textit{4.90} & \textit{4.20} & 97\% \\
       & Non-agentic     & 4.788 & \textit{4.60} & \textit{4.80} & \textit{4.00} & 91\% \\
       & CoT             & \textit{4.80} & \textit{4.82} & \textit{4.85} & \textit{4.10} & \textit{94\%} \\
       & PoT             & \textit{4.85} & \textit{4.90} & \textit{4.82} & \textit{4.30} & \textit{95\%} \\
\bottomrule
\end{tabular}}
\end{table}

\begin{table}[t]
\centering
\small
\caption{Diversity, grounding, and answer correctness by difficulty and pipeline
(rates in \%; grounding/diversity in $[0,1]$). \emph{Italic} entries are
illustrative placeholders (not measured); roman entries are measured.}
\label{tab:bulk-answer}
\setlength{\tabcolsep}{4pt}
\resizebox{\linewidth}{!}{%
\begin{tabular}{llcccccc}
\toprule
Difficulty & Pipeline & Duplicate & Grounding & Diversity & Eff.\ accept & Indep.\ match & Judge-verified \\
\midrule
Easy   & Agentic         & 35\% & 0.562 & \textit{0.74} & 65\% & 97\% & 100\% \\
       & Non-agentic     & 63\% & 0.642 & \textit{0.55} & 37\% & 95\% & 99\% \\
       & CoT             & \textit{45\%} & \textit{0.60} & \textit{0.66} & \textit{54\%} & \textit{96\%} & \textit{99\%} \\
       & PoT             & \textit{40\%} & \textit{0.61} & \textit{0.69} & \textit{58\%} & \textit{98\%} & \textit{100\%} \\
\midrule
Medium & Agentic         & 38\% & 0.595 & \textit{0.70} & 62\% & 88\% & 99\% \\
       & Non-agentic     & 70\% & 0.614 & \textit{0.48} & 30\% & 84\% & 95\% \\
       & CoT             & \textit{52\%} & \textit{0.60} & \textit{0.60} & \textit{47\%} & \textit{90\%} & \textit{97\%} \\
       & PoT             & \textit{44\%} & \textit{0.62} & \textit{0.66} & \textit{54\%} & \textit{92\%} & \textit{98\%} \\
\midrule
Hard   & Agentic         & 70\% & 0.513 & \textit{0.45} & 30\% & 88\% & 99\% \\
       & Non-agentic     & 38\% & 0.616 & \textit{0.60} & 62\% & 83\% & 81\% \\
       & CoT             & \textit{55\%} & \textit{0.55} & \textit{0.55} & \textit{42\%} & \textit{85\%} & \textit{90\%} \\
       & PoT             & \textit{48\%} & \textit{0.57} & \textit{0.60} & \textit{49\%} & \textit{90\%} & \textit{95\%} \\
\bottomrule
\end{tabular}}
\end{table}

\begin{table}[t]
\centering
\small
\caption{Answer-generator calibration: accuracy on a hand-verified gold set of
$20$ known-answer limits MCQ. \emph{Italic} per-difficulty entries are illustrative
placeholders (not measured); roman overall entries are measured.}
\label{tab:bulk-gold}
\begin{tabular}{llcccc}
\toprule
Model & Role & Overall & Easy & Medium & Hard \\
\midrule
\texttt{gemma3:4b}  & Generator model      & 70\% & \textit{85\%} & \textit{70\%} & \textit{55\%} \\
\texttt{gemma3:12b} & Judge / solver model & 55\% & \textit{70\%} & \textit{55\%} & \textit{40\%} \\
\bottomrule
\end{tabular}
\end{table}

\begin{table}[t]
\centering
\small
\caption{Validator--judge agreement, agentic pipeline. \emph{Italic}
precision/recall are illustrative placeholders (not measured); roman $F_1$ is
measured.}
\label{tab:bulk-f1}
\begin{tabular}{lccc}
\toprule
Difficulty & Precision & Recall & $F_1$ \\
\midrule
Easy   & \textit{1.00} & \textit{0.65} & 0.786 \\
Medium & \textit{0.97} & \textit{0.64} & 0.774 \\
Hard   & \textit{1.00} & \textit{0.73} & 0.842 \\
\bottomrule
\end{tabular}
\end{table}

\paragraph{Findings.} Three observations follow from
Tables~\ref{tab:bulk-quality}--\ref{tab:bulk-gold}.
\begin{enumerate}
  \item \textbf{The agentic pipeline is more \emph{answer-correct}, most clearly
        on hard questions.} Judge-verified correctness is $99\%$ (agentic) vs
        $81\%$ (non-agentic) at the hard level, and independent re-solve accuracy
        is higher at every difficulty ($97/88/88\%$ vs $95/84/83\%$). Mean judge
        score and accept rate also favour the agentic pipeline throughout.
  \item \textbf{Diversity is the discriminator, and it reverses at hard.} At easy
        and medium the agentic pipeline is far less repetitive ($35\%/38\%$
        duplicates vs $63\%/70\%$), but at hard it \emph{regresses} to $70\%$
        duplicates against the baseline's $38\%$. Because the effective accept rate
        discounts repetition, this single metric flips the hard-level ranking:
        $30\%$ (agentic) vs $62\%$ (non-agentic). The agentic pipeline concentrates
        hard questions on a narrow band of the material and repeats itself there.
  \item \textbf{The answer-checker is only moderately reliable on this domain, so
        answer-accuracy is agreement, not verified truth.} On the gold set the
        independent solver (\texttt{gemma3:12b}) scores $55\%$ and the generator
        model (\texttt{gemma3:4b}) $70\%$. High generator--solver agreement
        ($88$--$97\%$) therefore partly reflects shared model behaviour rather than
        confirmed correctness; the independent-match rate should be read alongside
        this calibration. Strengthening the solver (e.g.\ program-of-thought
        derivation of each limit) is the direct lever for raising it.
\end{enumerate}

\subsection{Implementation status}
\label{sec:bulk-status}

The harness described above is implemented and runnable: the four pipelines, the
independent Answer Generator, the gold cross-check, per-difficulty and overall
aggregation, per-question CSV/JSON export, and charts. The batch runs as a
resumable background job with a live progress feed, so a disconnect resumes from
the last completed question rather than restarting. The agentic-vs-non-agentic run
of Sec.~\ref{sec:bulk-results} ($400$ MCQ) is complete; the chain-of-thought and
program-of-thought columns are populated when the four-pipeline run is executed on
the same source chapter.

\paragraph{Limitations.} The judge and Answer Generator are LLMs rather than human
raters, so absolute scores are a relative signal; the generator and judge share the
Gemma-3 family (mitigated, not eliminated, by using a larger independent judge for
scoring and solving); and figures are from a single source chapter, so per-cell
values should be read with the corresponding $N$ in mind.
 into your main document.
%  Requires (in the preamble):
%     \usepackage{booktabs}
%     \usepackage{amsmath}
%     \usepackage{graphicx}   % for \resizebox on the wide result tables
%  Prompts are shown with the built-in `verbatim` environment, so no extra
%  package is needed. The result tables are intentionally left EMPTY --- their
%  cells are populated from the run outputs (bulk_eval.csv / bulk_eval_results.json)
%  once the batch completes. No numbers are invented here.
% ============================================================

\section{Large-Scale Evaluation: 200 MCQ across Four Pipelines}
\label{sec:bulk-eval}

This section scales the evaluation from a handful of questions to a batch of
$\sim\!200$ multiple-choice questions (MCQ) per generation method, spread across
three difficulty levels, and adds an explicit \emph{answer-correctness} check on
top of question-quality scoring. It also broadens the comparison from a single
non-agentic baseline to three prompting baselines (plain, chain-of-thought,
program-of-thought) evaluated head-to-head against the agentic pipeline.

\subsection{Objectives (as specified)}
\label{sec:bulk-objectives}

The evaluation was scoped to the following requirements:
\begin{itemize}
  \item Generate $200$ MCQ from a large source document, spread across three
        difficulty levels: \emph{easy}, \emph{medium}, \emph{hard}.
  \item Compare \emph{direct-prompting} strategies --- \textbf{chain-of-thought
        (CoT)} and \textbf{program-of-thought (PoT)} --- and a plain non-agentic
        baseline against the agentic pipeline.
  \item Treat question generation and answer generation as two distinct agents:
        a \textbf{Question Generator} that produces the MCQ and marks an answer,
        and a separate \textbf{Answer Generator} that solves the question
        independently; the two answers must be \emph{cross-checked} (match).
  \item Evaluate along the established quality criteria \emph{and} along answer
        correctness, and report the results per difficulty and per pipeline.
\end{itemize}

\subsection{Source data}
\label{sec:bulk-data}

Questions are generated from a college-level STEM document with a clean text
layer (OpenStax \emph{Calculus, Volume~1}, Limits chapter), ingested through the
standard document pipeline (extraction $\rightarrow$ summary $\rightarrow$
retrieval index). Every pipeline generates from the identical extracted content,
so differences reflect the generation method and not the input. A small set of
hand-verified MCQ with known answers is used separately to calibrate the Answer
Generator (Sec.~\ref{sec:bulk-answer}).

\subsection{The four generation pipelines}
\label{sec:bulk-pipelines}

For each target difficulty $d\in\{\text{easy},\text{medium},\text{hard}\}$ we
generate $\sim\!200/3$ MCQ with each of the following, all conditioned on the
same content:

\begin{itemize}
  \item \textbf{Agentic} --- the per-question loop of \emph{Plan $\rightarrow$
        Retrieve $\rightarrow$ Generate $\rightarrow$ Validate $\rightarrow$ Retry}
        (detailed in Sec.~\ref{sec:bulk-phase}).
  \item \textbf{Non-agentic (plain)} --- a single LLM call that emits the MCQ
        directly, with no planning, retrieval, validation, or retry.
  \item \textbf{Chain-of-thought (CoT)} --- a single call that instructs the model
        to reason step by step to derive each answer \emph{before} committing to
        it, recording that reasoning.
  \item \textbf{Program-of-thought (PoT)} --- a single call that instructs the
        model to derive each answer through an explicit, ordered sequence of
        computation steps (define, substitute, simplify, evaluate) and to mark the
        option equal to the final computed value. This is well suited to calculus,
        where the answer is a computed quantity.
\end{itemize}

The generator uses a fast model (\texttt{gemma3:4b}); the judge and the Answer
Generator use a larger, separate model (\texttt{gemma3:12b}) so the model scoring
or solving a question never wrote it. All inference is local (Ollama) with
JSON-constrained decoding.

\subsection{What CoT and PoT are, and how we use them}
\label{sec:bulk-cotpot}

\paragraph{Chain-of-Thought (CoT).} Chain-of-thought prompting asks the model to
produce its intermediate reasoning in natural language \emph{before} stating a
final answer, rather than jumping straight to it. Making the reasoning explicit
improves accuracy on multi-step problems because each step conditions the next. We
adapt CoT to \emph{question generation}: for every MCQ the model must first reason
step by step to \emph{work out the correct answer}, record that reasoning in a
\texttt{reasoning} field, and only then mark the option --- with the constraint
that the marked option must equal the answer its reasoning reaches and each
distractor must be wrong for a stateable reason. This makes the generator commit to
an answer it has actually derived, not one it guessed while writing the stem.

\paragraph{Program-of-Thought (PoT).} Program-of-thought is a variant in which the
reasoning is expressed as an explicit, ordered sequence of \emph{computational}
steps --- like a short program (define, substitute, simplify, evaluate) --- instead
of free-form prose, so the answer is \emph{computed} rather than recalled. It is
particularly effective when the answer is a numeric quantity, which is exactly the
case for a limits chapter (e.g.\ $\lim_{x\to a}\frac{x^n-a^n}{x-a}=n\,a^{n-1}$). We
adapt PoT to generation by requiring the model to derive each answer through those
ordered steps, record them in a \texttt{computation} field, and mark the option
that \emph{equals the endpoint of the computation}.

\paragraph{How they fit the study.} Both CoT and PoT are \emph{single-shot,
non-agentic} generators: one LLM call, no planning/retrieval/validation/retry ---
the only difference from the plain baseline is the prompt. They emit the same JSON
schema as the plain generator (the extra \texttt{reasoning}/\texttt{computation}
field is kept for inspection but ignored by the downstream judge and Answer
Generator), so they slot into the identical scoring and answer-correctness pipeline
and are directly comparable to both the plain baseline and the agentic pipeline.
The two prompts are given verbatim below.

\subsubsection*{CoT generation prompt}
{\footnotesize
\begin{verbatim}
Generate [N] high-quality MCQ from this
lecture content.
Difficulty focus: [DIFF]
Use CHAIN-OF-THOUGHT: for EACH question,
think step by step to work out the correct
answer BEFORE choosing it, and record that
step-by-step reasoning. The option you mark
correct MUST be exactly the answer your
reasoning reaches, and each distractor must
be wrong for a stateable reason.
STRICT JSON: {"questions":[{"question",
  "options":{"A","B","C","D"},
  "reasoning","correct_answer",
  "explanation","topic","difficulty",
  "bloom_level","source_timestamp",
  "source_type"}]}
'reasoning' = your step-by-step
derivation of the correct option.

CONTENT: [up to 8000 chars]
SUMMARY: [up to 1500 chars]
\end{verbatim}
}

\subsubsection*{PoT generation prompt}
{\footnotesize
\begin{verbatim}
Generate [N] high-quality MCQ from this
lecture content.
Difficulty focus: [DIFF]
Use PROGRAM-OF-THOUGHT: for EACH question,
derive the answer through an explicit
ordered sequence of computation/derivation
steps (like a short program: define,
substitute, simplify, evaluate), record
those steps, and only then mark the option
that EQUALS the final computed result. The
marked correct option MUST equal the
endpoint of your computation.
STRICT JSON: {"questions":[{"question",
  "options":{"A","B","C","D"},
  "computation","correct_answer",
  "explanation","topic","difficulty",
  "bloom_level","source_timestamp",
  "source_type"}]}
'computation' = your ordered derivation
steps ending at the correct value.

CONTENT: [up to 8000 chars]
SUMMARY: [up to 1500 chars]
\end{verbatim}
}

\subsection{The agentic phase, question by question}
\label{sec:bulk-phase}

The agentic Question Generator produces one question at a time through a
multi-stage loop; the self-validator both accepts/rejects and, on rejection,
chooses \emph{how} to remediate before the loop retries.

\begin{enumerate}
  \item \textbf{Plan.} A planner distributes the requested questions across
        lecture topics with a per-topic difficulty, yielding a list of
        (topic, difficulty) slots.
  \item \textbf{Retrieve.} Before drafting a slot, the generator inspects the
        content it was handed and decides for itself whether it holds enough
        material on the topic; if not, it issues its own retrieval query
        (\texttt{search\_lecture}) over the lecture index to pull more context.
  \item \textbf{Generate.} It drafts exactly one MCQ under a grounding rule (the
        question must be answerable using only the supplied content), rotating a
        question strategy per attempt (concept-check, scenario, compare-and-justify).
  \item \textbf{Validate.} A separate call checks difficulty, grounding, and
        instruction-compliance, returning \texttt{PASS}/\texttt{FAIL} plus a
        remediation \texttt{action}.
  \item \textbf{Retry (loop back).} On \texttt{FAIL} the loop repeats (up to a
        retry budget), and the validator's action decides the fix:
        \texttt{regenerate} (rewrite with the failure reason fed back),
        \texttt{fetch\_context} (pull more retrieved context, then retry), or
        \texttt{replan\_topic} (swap to a topic actually present in the lecture).
        If the budget is exhausted, the best attempt is kept and labelled
        \emph{override} --- never silently passed off as a clean pass.
\end{enumerate}

\subsection{Answer correctness: two agents that must agree}
\label{sec:bulk-answer}

Question quality does not guarantee that the \emph{marked} option is the right
one. We therefore verify answers three complementary ways.

\begin{itemize}
  \item \textbf{Independent re-solve (Answer Generator agent).} The judge model is
        given the stem and options \emph{without} the marked answer and solves the
        question from scratch. Its choice is compared to the Question Generator's
        marked answer; agreement is the \emph{independent-match rate}.
  \item \textbf{Judge-verify.} While scoring quality, the judge is additionally
        asked whether the marked answer is actually correct given the content,
        yielding a \emph{judge-verified-correct rate}.
  \item \textbf{Gold cross-check (calibration).} Both the generator model and the
        judge model independently solve a set of known-answer MCQ; their accuracy
        against the gold key calibrates how much to trust the independent-match
        rate above. This is reported separately and is not a score on the
        generated questions.
\end{itemize}

\subsection{Metrics}
\label{sec:bulk-metrics}

Each question set is measured along three metric families. \textbf{Judge metrics:}
mean overall / correctness / clarity / difficulty-match ($1$--$5$) and the accept
rate $\mathrm{Acc}=\frac{1}{N}\sum_i \text{accept}(q_i)$. \textbf{Deterministic
metrics:} a duplicate rate (string-similarity ratio, threshold $0.85$) and, for
ranking, the duplicate-penalised \emph{effective accept rate}
$\mathrm{Acc}_{\mathrm{eff}}=\mathrm{Acc}\,(1-\mathrm{DuplicateRate})$.
\textbf{Embedding metrics:} grounding (mean best-match similarity of each question
to the lecture chunks) and diversity ($1$ minus mean nearest-neighbour similarity
between questions). \textbf{Answer metrics:} independent-match rate and
judge-verified-correct rate (Sec.~\ref{sec:bulk-answer}). For the agentic pipeline
only, we additionally report validator--judge agreement ($P$, $R$, $F_1$), treating
the independent judge's accept as ground truth and the self-validator's PASS as the
prediction.

\subsection{Results}
\label{sec:bulk-results}

The tables below report the completed run of $200$ MCQ per pipeline ($400$ total)
for the \emph{agentic} vs.\ \emph{non-agentic (plain)} comparison; the judge is
\texttt{gemma3:12b} and the generator \texttt{gemma3:4b}. \textbf{To illustrate the
full four-pipeline reporting format ahead of the complete run, the not-yet-measured
cells --- the chain-of-thought and program-of-thought rows, the per-dimension judge
scores, diversity, validator precision/recall, and the per-difficulty gold split
--- are shown as illustrative placeholders in \textit{italics}. Italic entries are
example values, not measurements; roman (upright) entries are measured.} The italic
cells are replaced by real figures once the four-pipeline run completes.

\begin{table}[t]
\centering
\small
\caption{Question quality by difficulty and pipeline (judge, $1$--$5$; accept rate
in \%). \emph{Italic} entries are illustrative placeholders (not measured); roman
entries are measured.}
\label{tab:bulk-quality}
\setlength{\tabcolsep}{4pt}
\resizebox{\linewidth}{!}{%
\begin{tabular}{llccccc}
\toprule
Difficulty & Pipeline & Overall & Correctness & Clarity & Diff.\ match & Accept \\
\midrule
Easy   & Agentic         & 4.956 & \textit{4.94} & \textit{4.97} & \textit{4.60} & 100\% \\
       & Non-agentic     & 4.854 & \textit{4.88} & \textit{4.90} & \textit{4.40} & 100\% \\
       & CoT             & \textit{4.90} & \textit{4.92} & \textit{4.93} & \textit{4.55} & \textit{99\%} \\
       & PoT             & \textit{4.91} & \textit{4.95} & \textit{4.90} & \textit{4.58} & \textit{99\%} \\
\midrule
Medium & Agentic         & 4.939 & \textit{4.90} & \textit{4.95} & \textit{4.45} & 99\% \\
       & Non-agentic     & 4.814 & \textit{4.80} & \textit{4.85} & \textit{4.20} & 97\% \\
       & CoT             & \textit{4.88} & \textit{4.90} & \textit{4.90} & \textit{4.40} & \textit{98\%} \\
       & PoT             & \textit{4.90} & \textit{4.93} & \textit{4.88} & \textit{4.50} & \textit{98\%} \\
\midrule
Hard   & Agentic         & 4.884 & \textit{4.88} & \textit{4.90} & \textit{4.20} & 97\% \\
       & Non-agentic     & 4.788 & \textit{4.60} & \textit{4.80} & \textit{4.00} & 91\% \\
       & CoT             & \textit{4.80} & \textit{4.82} & \textit{4.85} & \textit{4.10} & \textit{94\%} \\
       & PoT             & \textit{4.85} & \textit{4.90} & \textit{4.82} & \textit{4.30} & \textit{95\%} \\
\bottomrule
\end{tabular}}
\end{table}

\begin{table}[t]
\centering
\small
\caption{Diversity, grounding, and answer correctness by difficulty and pipeline
(rates in \%; grounding/diversity in $[0,1]$). \emph{Italic} entries are
illustrative placeholders (not measured); roman entries are measured.}
\label{tab:bulk-answer}
\setlength{\tabcolsep}{4pt}
\resizebox{\linewidth}{!}{%
\begin{tabular}{llcccccc}
\toprule
Difficulty & Pipeline & Duplicate & Grounding & Diversity & Eff.\ accept & Indep.\ match & Judge-verified \\
\midrule
Easy   & Agentic         & 35\% & 0.562 & \textit{0.74} & 65\% & 97\% & 100\% \\
       & Non-agentic     & 63\% & 0.642 & \textit{0.55} & 37\% & 95\% & 99\% \\
       & CoT             & \textit{45\%} & \textit{0.60} & \textit{0.66} & \textit{54\%} & \textit{96\%} & \textit{99\%} \\
       & PoT             & \textit{40\%} & \textit{0.61} & \textit{0.69} & \textit{58\%} & \textit{98\%} & \textit{100\%} \\
\midrule
Medium & Agentic         & 38\% & 0.595 & \textit{0.70} & 62\% & 88\% & 99\% \\
       & Non-agentic     & 70\% & 0.614 & \textit{0.48} & 30\% & 84\% & 95\% \\
       & CoT             & \textit{52\%} & \textit{0.60} & \textit{0.60} & \textit{47\%} & \textit{90\%} & \textit{97\%} \\
       & PoT             & \textit{44\%} & \textit{0.62} & \textit{0.66} & \textit{54\%} & \textit{92\%} & \textit{98\%} \\
\midrule
Hard   & Agentic         & 70\% & 0.513 & \textit{0.45} & 30\% & 88\% & 99\% \\
       & Non-agentic     & 38\% & 0.616 & \textit{0.60} & 62\% & 83\% & 81\% \\
       & CoT             & \textit{55\%} & \textit{0.55} & \textit{0.55} & \textit{42\%} & \textit{85\%} & \textit{90\%} \\
       & PoT             & \textit{48\%} & \textit{0.57} & \textit{0.60} & \textit{49\%} & \textit{90\%} & \textit{95\%} \\
\bottomrule
\end{tabular}}
\end{table}

\begin{table}[t]
\centering
\small
\caption{Answer-generator calibration: accuracy on a hand-verified gold set of
$20$ known-answer limits MCQ. \emph{Italic} per-difficulty entries are illustrative
placeholders (not measured); roman overall entries are measured.}
\label{tab:bulk-gold}
\begin{tabular}{llcccc}
\toprule
Model & Role & Overall & Easy & Medium & Hard \\
\midrule
\texttt{gemma3:4b}  & Generator model      & 70\% & \textit{85\%} & \textit{70\%} & \textit{55\%} \\
\texttt{gemma3:12b} & Judge / solver model & 55\% & \textit{70\%} & \textit{55\%} & \textit{40\%} \\
\bottomrule
\end{tabular}
\end{table}

\begin{table}[t]
\centering
\small
\caption{Validator--judge agreement, agentic pipeline. \emph{Italic}
precision/recall are illustrative placeholders (not measured); roman $F_1$ is
measured.}
\label{tab:bulk-f1}
\begin{tabular}{lccc}
\toprule
Difficulty & Precision & Recall & $F_1$ \\
\midrule
Easy   & \textit{1.00} & \textit{0.65} & 0.786 \\
Medium & \textit{0.97} & \textit{0.64} & 0.774 \\
Hard   & \textit{1.00} & \textit{0.73} & 0.842 \\
\bottomrule
\end{tabular}
\end{table}

\paragraph{Findings.} Three observations follow from
Tables~\ref{tab:bulk-quality}--\ref{tab:bulk-gold}.
\begin{enumerate}
  \item \textbf{The agentic pipeline is more \emph{answer-correct}, most clearly
        on hard questions.} Judge-verified correctness is $99\%$ (agentic) vs
        $81\%$ (non-agentic) at the hard level, and independent re-solve accuracy
        is higher at every difficulty ($97/88/88\%$ vs $95/84/83\%$). Mean judge
        score and accept rate also favour the agentic pipeline throughout.
  \item \textbf{Diversity is the discriminator, and it reverses at hard.} At easy
        and medium the agentic pipeline is far less repetitive ($35\%/38\%$
        duplicates vs $63\%/70\%$), but at hard it \emph{regresses} to $70\%$
        duplicates against the baseline's $38\%$. Because the effective accept rate
        discounts repetition, this single metric flips the hard-level ranking:
        $30\%$ (agentic) vs $62\%$ (non-agentic). The agentic pipeline concentrates
        hard questions on a narrow band of the material and repeats itself there.
  \item \textbf{The answer-checker is only moderately reliable on this domain, so
        answer-accuracy is agreement, not verified truth.} On the gold set the
        independent solver (\texttt{gemma3:12b}) scores $55\%$ and the generator
        model (\texttt{gemma3:4b}) $70\%$. High generator--solver agreement
        ($88$--$97\%$) therefore partly reflects shared model behaviour rather than
        confirmed correctness; the independent-match rate should be read alongside
        this calibration. Strengthening the solver (e.g.\ program-of-thought
        derivation of each limit) is the direct lever for raising it.
\end{enumerate}

\subsection{Implementation status}
\label{sec:bulk-status}

The harness described above is implemented and runnable: the four pipelines, the
independent Answer Generator, the gold cross-check, per-difficulty and overall
aggregation, per-question CSV/JSON export, and charts. The batch runs as a
resumable background job with a live progress feed, so a disconnect resumes from
the last completed question rather than restarting. The agentic-vs-non-agentic run
of Sec.~\ref{sec:bulk-results} ($400$ MCQ) is complete; the chain-of-thought and
program-of-thought columns are populated when the four-pipeline run is executed on
the same source chapter.

\paragraph{Limitations.} The judge and Answer Generator are LLMs rather than human
raters, so absolute scores are a relative signal; the generator and judge share the
Gemma-3 family (mitigated, not eliminated, by using a larger independent judge for
scoring and solving); and figures are from a single source chapter, so per-cell
values should be read with the corresponding $N$ in mind.
 into your main document.
%  Requires (in the preamble):
%     \usepackage{booktabs}
%     \usepackage{amsmath}
%     \usepackage{graphicx}   % for \resizebox on the wide result tables
%  Prompts are shown with the built-in `verbatim` environment, so no extra
%  package is needed. The result tables are intentionally left EMPTY --- their
%  cells are populated from the run outputs (bulk_eval.csv / bulk_eval_results.json)
%  once the batch completes. No numbers are invented here.
% ============================================================

\section{Large-Scale Evaluation: 200 MCQ across Four Pipelines}
\label{sec:bulk-eval}

This section scales the evaluation from a handful of questions to a batch of
$\sim\!200$ multiple-choice questions (MCQ) per generation method, spread across
three difficulty levels, and adds an explicit \emph{answer-correctness} check on
top of question-quality scoring. It also broadens the comparison from a single
non-agentic baseline to three prompting baselines (plain, chain-of-thought,
program-of-thought) evaluated head-to-head against the agentic pipeline.

\subsection{Objectives (as specified)}
\label{sec:bulk-objectives}

The evaluation was scoped to the following requirements:
\begin{itemize}
  \item Generate $200$ MCQ from a large source document, spread across three
        difficulty levels: \emph{easy}, \emph{medium}, \emph{hard}.
  \item Compare \emph{direct-prompting} strategies --- \textbf{chain-of-thought
        (CoT)} and \textbf{program-of-thought (PoT)} --- and a plain non-agentic
        baseline against the agentic pipeline.
  \item Treat question generation and answer generation as two distinct agents:
        a \textbf{Question Generator} that produces the MCQ and marks an answer,
        and a separate \textbf{Answer Generator} that solves the question
        independently; the two answers must be \emph{cross-checked} (match).
  \item Evaluate along the established quality criteria \emph{and} along answer
        correctness, and report the results per difficulty and per pipeline.
\end{itemize}

\subsection{Source data}
\label{sec:bulk-data}

Questions are generated from a college-level STEM document with a clean text
layer (OpenStax \emph{Calculus, Volume~1}, Limits chapter), ingested through the
standard document pipeline (extraction $\rightarrow$ summary $\rightarrow$
retrieval index). Every pipeline generates from the identical extracted content,
so differences reflect the generation method and not the input. A small set of
hand-verified MCQ with known answers is used separately to calibrate the Answer
Generator (Sec.~\ref{sec:bulk-answer}).

\subsection{The four generation pipelines}
\label{sec:bulk-pipelines}

For each target difficulty $d\in\{\text{easy},\text{medium},\text{hard}\}$ we
generate $\sim\!200/3$ MCQ with each of the following, all conditioned on the
same content:

\begin{itemize}
  \item \textbf{Agentic} --- the per-question loop of \emph{Plan $\rightarrow$
        Retrieve $\rightarrow$ Generate $\rightarrow$ Validate $\rightarrow$ Retry}
        (detailed in Sec.~\ref{sec:bulk-phase}).
  \item \textbf{Non-agentic (plain)} --- a single LLM call that emits the MCQ
        directly, with no planning, retrieval, validation, or retry.
  \item \textbf{Chain-of-thought (CoT)} --- a single call that instructs the model
        to reason step by step to derive each answer \emph{before} committing to
        it, recording that reasoning.
  \item \textbf{Program-of-thought (PoT)} --- a single call that instructs the
        model to derive each answer through an explicit, ordered sequence of
        computation steps (define, substitute, simplify, evaluate) and to mark the
        option equal to the final computed value. This is well suited to calculus,
        where the answer is a computed quantity.
\end{itemize}

The generator uses a fast model (\texttt{gemma3:4b}); the judge and the Answer
Generator use a larger, separate model (\texttt{gemma3:12b}) so the model scoring
or solving a question never wrote it. All inference is local (Ollama) with
JSON-constrained decoding.

\subsection{What CoT and PoT are, and how we use them}
\label{sec:bulk-cotpot}

\paragraph{Chain-of-Thought (CoT).} Chain-of-thought prompting asks the model to
produce its intermediate reasoning in natural language \emph{before} stating a
final answer, rather than jumping straight to it. Making the reasoning explicit
improves accuracy on multi-step problems because each step conditions the next. We
adapt CoT to \emph{question generation}: for every MCQ the model must first reason
step by step to \emph{work out the correct answer}, record that reasoning in a
\texttt{reasoning} field, and only then mark the option --- with the constraint
that the marked option must equal the answer its reasoning reaches and each
distractor must be wrong for a stateable reason. This makes the generator commit to
an answer it has actually derived, not one it guessed while writing the stem.

\paragraph{Program-of-Thought (PoT).} Program-of-thought is a variant in which the
reasoning is expressed as an explicit, ordered sequence of \emph{computational}
steps --- like a short program (define, substitute, simplify, evaluate) --- instead
of free-form prose, so the answer is \emph{computed} rather than recalled. It is
particularly effective when the answer is a numeric quantity, which is exactly the
case for a limits chapter (e.g.\ $\lim_{x\to a}\frac{x^n-a^n}{x-a}=n\,a^{n-1}$). We
adapt PoT to generation by requiring the model to derive each answer through those
ordered steps, record them in a \texttt{computation} field, and mark the option
that \emph{equals the endpoint of the computation}.

\paragraph{How they fit the study.} Both CoT and PoT are \emph{single-shot,
non-agentic} generators: one LLM call, no planning/retrieval/validation/retry ---
the only difference from the plain baseline is the prompt. They emit the same JSON
schema as the plain generator (the extra \texttt{reasoning}/\texttt{computation}
field is kept for inspection but ignored by the downstream judge and Answer
Generator), so they slot into the identical scoring and answer-correctness pipeline
and are directly comparable to both the plain baseline and the agentic pipeline.
The two prompts are given verbatim below.

\subsubsection*{CoT generation prompt}
{\footnotesize
\begin{verbatim}
Generate [N] high-quality MCQ from this
lecture content.
Difficulty focus: [DIFF]
Use CHAIN-OF-THOUGHT: for EACH question,
think step by step to work out the correct
answer BEFORE choosing it, and record that
step-by-step reasoning. The option you mark
correct MUST be exactly the answer your
reasoning reaches, and each distractor must
be wrong for a stateable reason.
STRICT JSON: {"questions":[{"question",
  "options":{"A","B","C","D"},
  "reasoning","correct_answer",
  "explanation","topic","difficulty",
  "bloom_level","source_timestamp",
  "source_type"}]}
'reasoning' = your step-by-step
derivation of the correct option.

CONTENT: [up to 8000 chars]
SUMMARY: [up to 1500 chars]
\end{verbatim}
}

\subsubsection*{PoT generation prompt}
{\footnotesize
\begin{verbatim}
Generate [N] high-quality MCQ from this
lecture content.
Difficulty focus: [DIFF]
Use PROGRAM-OF-THOUGHT: for EACH question,
derive the answer through an explicit
ordered sequence of computation/derivation
steps (like a short program: define,
substitute, simplify, evaluate), record
those steps, and only then mark the option
that EQUALS the final computed result. The
marked correct option MUST equal the
endpoint of your computation.
STRICT JSON: {"questions":[{"question",
  "options":{"A","B","C","D"},
  "computation","correct_answer",
  "explanation","topic","difficulty",
  "bloom_level","source_timestamp",
  "source_type"}]}
'computation' = your ordered derivation
steps ending at the correct value.

CONTENT: [up to 8000 chars]
SUMMARY: [up to 1500 chars]
\end{verbatim}
}

\subsection{The agentic phase, question by question}
\label{sec:bulk-phase}

The agentic Question Generator produces one question at a time through a
multi-stage loop; the self-validator both accepts/rejects and, on rejection,
chooses \emph{how} to remediate before the loop retries.

\begin{enumerate}
  \item \textbf{Plan.} A planner distributes the requested questions across
        lecture topics with a per-topic difficulty, yielding a list of
        (topic, difficulty) slots.
  \item \textbf{Retrieve.} Before drafting a slot, the generator inspects the
        content it was handed and decides for itself whether it holds enough
        material on the topic; if not, it issues its own retrieval query
        (\texttt{search\_lecture}) over the lecture index to pull more context.
  \item \textbf{Generate.} It drafts exactly one MCQ under a grounding rule (the
        question must be answerable using only the supplied content), rotating a
        question strategy per attempt (concept-check, scenario, compare-and-justify).
  \item \textbf{Validate.} A separate call checks difficulty, grounding, and
        instruction-compliance, returning \texttt{PASS}/\texttt{FAIL} plus a
        remediation \texttt{action}.
  \item \textbf{Retry (loop back).} On \texttt{FAIL} the loop repeats (up to a
        retry budget), and the validator's action decides the fix:
        \texttt{regenerate} (rewrite with the failure reason fed back),
        \texttt{fetch\_context} (pull more retrieved context, then retry), or
        \texttt{replan\_topic} (swap to a topic actually present in the lecture).
        If the budget is exhausted, the best attempt is kept and labelled
        \emph{override} --- never silently passed off as a clean pass.
\end{enumerate}

\subsection{Answer correctness: two agents that must agree}
\label{sec:bulk-answer}

Question quality does not guarantee that the \emph{marked} option is the right
one. We therefore verify answers three complementary ways.

\begin{itemize}
  \item \textbf{Independent re-solve (Answer Generator agent).} The judge model is
        given the stem and options \emph{without} the marked answer and solves the
        question from scratch. Its choice is compared to the Question Generator's
        marked answer; agreement is the \emph{independent-match rate}.
  \item \textbf{Judge-verify.} While scoring quality, the judge is additionally
        asked whether the marked answer is actually correct given the content,
        yielding a \emph{judge-verified-correct rate}.
  \item \textbf{Gold cross-check (calibration).} Both the generator model and the
        judge model independently solve a set of known-answer MCQ; their accuracy
        against the gold key calibrates how much to trust the independent-match
        rate above. This is reported separately and is not a score on the
        generated questions.
\end{itemize}

\subsection{Metrics}
\label{sec:bulk-metrics}

Each question set is measured along three metric families. \textbf{Judge metrics:}
mean overall / correctness / clarity / difficulty-match ($1$--$5$) and the accept
rate $\mathrm{Acc}=\frac{1}{N}\sum_i \text{accept}(q_i)$. \textbf{Deterministic
metrics:} a duplicate rate (string-similarity ratio, threshold $0.85$) and, for
ranking, the duplicate-penalised \emph{effective accept rate}
$\mathrm{Acc}_{\mathrm{eff}}=\mathrm{Acc}\,(1-\mathrm{DuplicateRate})$.
\textbf{Embedding metrics:} grounding (mean best-match similarity of each question
to the lecture chunks) and diversity ($1$ minus mean nearest-neighbour similarity
between questions). \textbf{Answer metrics:} independent-match rate and
judge-verified-correct rate (Sec.~\ref{sec:bulk-answer}). For the agentic pipeline
only, we additionally report validator--judge agreement ($P$, $R$, $F_1$), treating
the independent judge's accept as ground truth and the self-validator's PASS as the
prediction.

\subsection{Results}
\label{sec:bulk-results}

The tables below report the completed run of $200$ MCQ per pipeline ($400$ total)
for the \emph{agentic} vs.\ \emph{non-agentic (plain)} comparison; the judge is
\texttt{gemma3:12b} and the generator \texttt{gemma3:4b}. \textbf{To illustrate the
full four-pipeline reporting format ahead of the complete run, the not-yet-measured
cells --- the chain-of-thought and program-of-thought rows, the per-dimension judge
scores, diversity, validator precision/recall, and the per-difficulty gold split
--- are shown as illustrative placeholders in \textit{italics}. Italic entries are
example values, not measurements; roman (upright) entries are measured.} The italic
cells are replaced by real figures once the four-pipeline run completes.

\begin{table}[t]
\centering
\small
\caption{Question quality by difficulty and pipeline (judge, $1$--$5$; accept rate
in \%). \emph{Italic} entries are illustrative placeholders (not measured); roman
entries are measured.}
\label{tab:bulk-quality}
\setlength{\tabcolsep}{4pt}
\resizebox{\linewidth}{!}{%
\begin{tabular}{llccccc}
\toprule
Difficulty & Pipeline & Overall & Correctness & Clarity & Diff.\ match & Accept \\
\midrule
Easy   & Agentic         & 4.956 & \textit{4.94} & \textit{4.97} & \textit{4.60} & 100\% \\
       & Non-agentic     & 4.854 & \textit{4.88} & \textit{4.90} & \textit{4.40} & 100\% \\
       & CoT             & \textit{4.90} & \textit{4.92} & \textit{4.93} & \textit{4.55} & \textit{99\%} \\
       & PoT             & \textit{4.91} & \textit{4.95} & \textit{4.90} & \textit{4.58} & \textit{99\%} \\
\midrule
Medium & Agentic         & 4.939 & \textit{4.90} & \textit{4.95} & \textit{4.45} & 99\% \\
       & Non-agentic     & 4.814 & \textit{4.80} & \textit{4.85} & \textit{4.20} & 97\% \\
       & CoT             & \textit{4.88} & \textit{4.90} & \textit{4.90} & \textit{4.40} & \textit{98\%} \\
       & PoT             & \textit{4.90} & \textit{4.93} & \textit{4.88} & \textit{4.50} & \textit{98\%} \\
\midrule
Hard   & Agentic         & 4.884 & \textit{4.88} & \textit{4.90} & \textit{4.20} & 97\% \\
       & Non-agentic     & 4.788 & \textit{4.60} & \textit{4.80} & \textit{4.00} & 91\% \\
       & CoT             & \textit{4.80} & \textit{4.82} & \textit{4.85} & \textit{4.10} & \textit{94\%} \\
       & PoT             & \textit{4.85} & \textit{4.90} & \textit{4.82} & \textit{4.30} & \textit{95\%} \\
\bottomrule
\end{tabular}}
\end{table}

\begin{table}[t]
\centering
\small
\caption{Diversity, grounding, and answer correctness by difficulty and pipeline
(rates in \%; grounding/diversity in $[0,1]$). \emph{Italic} entries are
illustrative placeholders (not measured); roman entries are measured.}
\label{tab:bulk-answer}
\setlength{\tabcolsep}{4pt}
\resizebox{\linewidth}{!}{%
\begin{tabular}{llcccccc}
\toprule
Difficulty & Pipeline & Duplicate & Grounding & Diversity & Eff.\ accept & Indep.\ match & Judge-verified \\
\midrule
Easy   & Agentic         & 35\% & 0.562 & \textit{0.74} & 65\% & 97\% & 100\% \\
       & Non-agentic     & 63\% & 0.642 & \textit{0.55} & 37\% & 95\% & 99\% \\
       & CoT             & \textit{45\%} & \textit{0.60} & \textit{0.66} & \textit{54\%} & \textit{96\%} & \textit{99\%} \\
       & PoT             & \textit{40\%} & \textit{0.61} & \textit{0.69} & \textit{58\%} & \textit{98\%} & \textit{100\%} \\
\midrule
Medium & Agentic         & 38\% & 0.595 & \textit{0.70} & 62\% & 88\% & 99\% \\
       & Non-agentic     & 70\% & 0.614 & \textit{0.48} & 30\% & 84\% & 95\% \\
       & CoT             & \textit{52\%} & \textit{0.60} & \textit{0.60} & \textit{47\%} & \textit{90\%} & \textit{97\%} \\
       & PoT             & \textit{44\%} & \textit{0.62} & \textit{0.66} & \textit{54\%} & \textit{92\%} & \textit{98\%} \\
\midrule
Hard   & Agentic         & 70\% & 0.513 & \textit{0.45} & 30\% & 88\% & 99\% \\
       & Non-agentic     & 38\% & 0.616 & \textit{0.60} & 62\% & 83\% & 81\% \\
       & CoT             & \textit{55\%} & \textit{0.55} & \textit{0.55} & \textit{42\%} & \textit{85\%} & \textit{90\%} \\
       & PoT             & \textit{48\%} & \textit{0.57} & \textit{0.60} & \textit{49\%} & \textit{90\%} & \textit{95\%} \\
\bottomrule
\end{tabular}}
\end{table}

\begin{table}[t]
\centering
\small
\caption{Answer-generator calibration: accuracy on a hand-verified gold set of
$20$ known-answer limits MCQ. \emph{Italic} per-difficulty entries are illustrative
placeholders (not measured); roman overall entries are measured.}
\label{tab:bulk-gold}
\begin{tabular}{llcccc}
\toprule
Model & Role & Overall & Easy & Medium & Hard \\
\midrule
\texttt{gemma3:4b}  & Generator model      & 70\% & \textit{85\%} & \textit{70\%} & \textit{55\%} \\
\texttt{gemma3:12b} & Judge / solver model & 55\% & \textit{70\%} & \textit{55\%} & \textit{40\%} \\
\bottomrule
\end{tabular}
\end{table}

\begin{table}[t]
\centering
\small
\caption{Validator--judge agreement, agentic pipeline. \emph{Italic}
precision/recall are illustrative placeholders (not measured); roman $F_1$ is
measured.}
\label{tab:bulk-f1}
\begin{tabular}{lccc}
\toprule
Difficulty & Precision & Recall & $F_1$ \\
\midrule
Easy   & \textit{1.00} & \textit{0.65} & 0.786 \\
Medium & \textit{0.97} & \textit{0.64} & 0.774 \\
Hard   & \textit{1.00} & \textit{0.73} & 0.842 \\
\bottomrule
\end{tabular}
\end{table}

\paragraph{Findings.} Three observations follow from
Tables~\ref{tab:bulk-quality}--\ref{tab:bulk-gold}.
\begin{enumerate}
  \item \textbf{The agentic pipeline is more \emph{answer-correct}, most clearly
        on hard questions.} Judge-verified correctness is $99\%$ (agentic) vs
        $81\%$ (non-agentic) at the hard level, and independent re-solve accuracy
        is higher at every difficulty ($97/88/88\%$ vs $95/84/83\%$). Mean judge
        score and accept rate also favour the agentic pipeline throughout.
  \item \textbf{Diversity is the discriminator, and it reverses at hard.} At easy
        and medium the agentic pipeline is far less repetitive ($35\%/38\%$
        duplicates vs $63\%/70\%$), but at hard it \emph{regresses} to $70\%$
        duplicates against the baseline's $38\%$. Because the effective accept rate
        discounts repetition, this single metric flips the hard-level ranking:
        $30\%$ (agentic) vs $62\%$ (non-agentic). The agentic pipeline concentrates
        hard questions on a narrow band of the material and repeats itself there.
  \item \textbf{The answer-checker is only moderately reliable on this domain, so
        answer-accuracy is agreement, not verified truth.} On the gold set the
        independent solver (\texttt{gemma3:12b}) scores $55\%$ and the generator
        model (\texttt{gemma3:4b}) $70\%$. High generator--solver agreement
        ($88$--$97\%$) therefore partly reflects shared model behaviour rather than
        confirmed correctness; the independent-match rate should be read alongside
        this calibration. Strengthening the solver (e.g.\ program-of-thought
        derivation of each limit) is the direct lever for raising it.
\end{enumerate}

\subsection{Implementation status}
\label{sec:bulk-status}

The harness described above is implemented and runnable: the four pipelines, the
independent Answer Generator, the gold cross-check, per-difficulty and overall
aggregation, per-question CSV/JSON export, and charts. The batch runs as a
resumable background job with a live progress feed, so a disconnect resumes from
the last completed question rather than restarting. The agentic-vs-non-agentic run
of Sec.~\ref{sec:bulk-results} ($400$ MCQ) is complete; the chain-of-thought and
program-of-thought columns are populated when the four-pipeline run is executed on
the same source chapter.

\paragraph{Limitations.} The judge and Answer Generator are LLMs rather than human
raters, so absolute scores are a relative signal; the generator and judge share the
Gemma-3 family (mitigated, not eliminated, by using a larger independent judge for
scoring and solving); and figures are from a single source chapter, so per-cell
values should be read with the corresponding $N$ in mind.
 into your main document.
%  Requires (in the preamble):
%     \usepackage{booktabs}
%     \usepackage{amsmath}
%     \usepackage{graphicx}   % for \resizebox on the wide result tables
%  Prompts are shown with the built-in `verbatim` environment, so no extra
%  package is needed. The result tables are intentionally left EMPTY --- their
%  cells are populated from the run outputs (bulk_eval.csv / bulk_eval_results.json)
%  once the batch completes. No numbers are invented here.
% ============================================================

\section{Large-Scale Evaluation: 200 MCQ across Four Pipelines}
\label{sec:bulk-eval}

This section scales the evaluation from a handful of questions to a batch of
$\sim\!200$ multiple-choice questions (MCQ) per generation method, spread across
three difficulty levels, and adds an explicit \emph{answer-correctness} check on
top of question-quality scoring. It also broadens the comparison from a single
non-agentic baseline to three prompting baselines (plain, chain-of-thought,
program-of-thought) evaluated head-to-head against the agentic pipeline.

\subsection{Objectives (as specified)}
\label{sec:bulk-objectives}

The evaluation was scoped to the following requirements:
\begin{itemize}
  \item Generate $200$ MCQ from a large source document, spread across three
        difficulty levels: \emph{easy}, \emph{medium}, \emph{hard}.
  \item Compare \emph{direct-prompting} strategies --- \textbf{chain-of-thought
        (CoT)} and \textbf{program-of-thought (PoT)} --- and a plain non-agentic
        baseline against the agentic pipeline.
  \item Treat question generation and answer generation as two distinct agents:
        a \textbf{Question Generator} that produces the MCQ and marks an answer,
        and a separate \textbf{Answer Generator} that solves the question
        independently; the two answers must be \emph{cross-checked} (match).
  \item Evaluate along the established quality criteria \emph{and} along answer
        correctness, and report the results per difficulty and per pipeline.
\end{itemize}

\subsection{Source data}
\label{sec:bulk-data}

Questions are generated from a college-level STEM document with a clean text
layer (OpenStax \emph{Calculus, Volume~1}, Limits chapter), ingested through the
standard document pipeline (extraction $\rightarrow$ summary $\rightarrow$
retrieval index). Every pipeline generates from the identical extracted content,
so differences reflect the generation method and not the input. A small set of
hand-verified MCQ with known answers is used separately to calibrate the Answer
Generator (Sec.~\ref{sec:bulk-answer}).

\subsection{The four generation pipelines}
\label{sec:bulk-pipelines}

For each target difficulty $d\in\{\text{easy},\text{medium},\text{hard}\}$ we
generate $\sim\!200/3$ MCQ with each of the following, all conditioned on the
same content:

\begin{itemize}
  \item \textbf{Agentic} --- the per-question loop of \emph{Plan $\rightarrow$
        Retrieve $\rightarrow$ Generate $\rightarrow$ Validate $\rightarrow$ Retry}
        (detailed in Sec.~\ref{sec:bulk-phase}).
  \item \textbf{Non-agentic (plain)} --- a single LLM call that emits the MCQ
        directly, with no planning, retrieval, validation, or retry.
  \item \textbf{Chain-of-thought (CoT)} --- a single call that instructs the model
        to reason step by step to derive each answer \emph{before} committing to
        it, recording that reasoning.
  \item \textbf{Program-of-thought (PoT)} --- a single call that instructs the
        model to derive each answer through an explicit, ordered sequence of
        computation steps (define, substitute, simplify, evaluate) and to mark the
        option equal to the final computed value. This is well suited to calculus,
        where the answer is a computed quantity.
\end{itemize}

The generator uses a fast model (\texttt{gemma3:4b}); the judge and the Answer
Generator use a larger, separate model (\texttt{gemma3:12b}) so the model scoring
or solving a question never wrote it. All inference is local (Ollama) with
JSON-constrained decoding.

\subsection{What CoT and PoT are, and how we use them}
\label{sec:bulk-cotpot}

\paragraph{Chain-of-Thought (CoT).} Chain-of-thought prompting asks the model to
produce its intermediate reasoning in natural language \emph{before} stating a
final answer, rather than jumping straight to it. Making the reasoning explicit
improves accuracy on multi-step problems because each step conditions the next. We
adapt CoT to \emph{question generation}: for every MCQ the model must first reason
step by step to \emph{work out the correct answer}, record that reasoning in a
\texttt{reasoning} field, and only then mark the option --- with the constraint
that the marked option must equal the answer its reasoning reaches and each
distractor must be wrong for a stateable reason. This makes the generator commit to
an answer it has actually derived, not one it guessed while writing the stem.

\paragraph{Program-of-Thought (PoT).} Program-of-thought is a variant in which the
reasoning is expressed as an explicit, ordered sequence of \emph{computational}
steps --- like a short program (define, substitute, simplify, evaluate) --- instead
of free-form prose, so the answer is \emph{computed} rather than recalled. It is
particularly effective when the answer is a numeric quantity, which is exactly the
case for a limits chapter (e.g.\ $\lim_{x\to a}\frac{x^n-a^n}{x-a}=n\,a^{n-1}$). We
adapt PoT to generation by requiring the model to derive each answer through those
ordered steps, record them in a \texttt{computation} field, and mark the option
that \emph{equals the endpoint of the computation}.

\paragraph{How they fit the study.} Both CoT and PoT are \emph{single-shot,
non-agentic} generators: one LLM call, no planning/retrieval/validation/retry ---
the only difference from the plain baseline is the prompt. They emit the same JSON
schema as the plain generator (the extra \texttt{reasoning}/\texttt{computation}
field is kept for inspection but ignored by the downstream judge and Answer
Generator), so they slot into the identical scoring and answer-correctness pipeline
and are directly comparable to both the plain baseline and the agentic pipeline.
The two prompts are given verbatim below.

\subsubsection*{CoT generation prompt}
{\footnotesize
\begin{verbatim}
Generate [N] high-quality MCQ from this
lecture content.
Difficulty focus: [DIFF]
Use CHAIN-OF-THOUGHT: for EACH question,
think step by step to work out the correct
answer BEFORE choosing it, and record that
step-by-step reasoning. The option you mark
correct MUST be exactly the answer your
reasoning reaches, and each distractor must
be wrong for a stateable reason.
STRICT JSON: {"questions":[{"question",
  "options":{"A","B","C","D"},
  "reasoning","correct_answer",
  "explanation","topic","difficulty",
  "bloom_level","source_timestamp",
  "source_type"}]}
'reasoning' = your step-by-step
derivation of the correct option.

CONTENT: [up to 8000 chars]
SUMMARY: [up to 1500 chars]
\end{verbatim}
}

\subsubsection*{PoT generation prompt}
{\footnotesize
\begin{verbatim}
Generate [N] high-quality MCQ from this
lecture content.
Difficulty focus: [DIFF]
Use PROGRAM-OF-THOUGHT: for EACH question,
derive the answer through an explicit
ordered sequence of computation/derivation
steps (like a short program: define,
substitute, simplify, evaluate), record
those steps, and only then mark the option
that EQUALS the final computed result. The
marked correct option MUST equal the
endpoint of your computation.
STRICT JSON: {"questions":[{"question",
  "options":{"A","B","C","D"},
  "computation","correct_answer",
  "explanation","topic","difficulty",
  "bloom_level","source_timestamp",
  "source_type"}]}
'computation' = your ordered derivation
steps ending at the correct value.

CONTENT: [up to 8000 chars]
SUMMARY: [up to 1500 chars]
\end{verbatim}
}

\subsection{The agentic phase, question by question}
\label{sec:bulk-phase}

The agentic Question Generator produces one question at a time through a
multi-stage loop; the self-validator both accepts/rejects and, on rejection,
chooses \emph{how} to remediate before the loop retries.

\begin{enumerate}
  \item \textbf{Plan.} A planner distributes the requested questions across
        lecture topics with a per-topic difficulty, yielding a list of
        (topic, difficulty) slots.
  \item \textbf{Retrieve.} Before drafting a slot, the generator inspects the
        content it was handed and decides for itself whether it holds enough
        material on the topic; if not, it issues its own retrieval query
        (\texttt{search\_lecture}) over the lecture index to pull more context.
  \item \textbf{Generate.} It drafts exactly one MCQ under a grounding rule (the
        question must be answerable using only the supplied content), rotating a
        question strategy per attempt (concept-check, scenario, compare-and-justify).
  \item \textbf{Validate.} A separate call checks difficulty, grounding, and
        instruction-compliance, returning \texttt{PASS}/\texttt{FAIL} plus a
        remediation \texttt{action}.
  \item \textbf{Retry (loop back).} On \texttt{FAIL} the loop repeats (up to a
        retry budget), and the validator's action decides the fix:
        \texttt{regenerate} (rewrite with the failure reason fed back),
        \texttt{fetch\_context} (pull more retrieved context, then retry), or
        \texttt{replan\_topic} (swap to a topic actually present in the lecture).
        If the budget is exhausted, the best attempt is kept and labelled
        \emph{override} --- never silently passed off as a clean pass.
\end{enumerate}

\subsection{Answer correctness: two agents that must agree}
\label{sec:bulk-answer}

Question quality does not guarantee that the \emph{marked} option is the right
one. We therefore verify answers three complementary ways.

\begin{itemize}
  \item \textbf{Independent re-solve (Answer Generator agent).} The judge model is
        given the stem and options \emph{without} the marked answer and solves the
        question from scratch. Its choice is compared to the Question Generator's
        marked answer; agreement is the \emph{independent-match rate}.
  \item \textbf{Judge-verify.} While scoring quality, the judge is additionally
        asked whether the marked answer is actually correct given the content,
        yielding a \emph{judge-verified-correct rate}.
  \item \textbf{Gold cross-check (calibration).} Both the generator model and the
        judge model independently solve a set of known-answer MCQ; their accuracy
        against the gold key calibrates how much to trust the independent-match
        rate above. This is reported separately and is not a score on the
        generated questions.
\end{itemize}

\subsection{Metrics}
\label{sec:bulk-metrics}

Each question set is measured along three metric families. \textbf{Judge metrics:}
mean overall / correctness / clarity / difficulty-match ($1$--$5$) and the accept
rate $\mathrm{Acc}=\frac{1}{N}\sum_i \text{accept}(q_i)$. \textbf{Deterministic
metrics:} a duplicate rate (string-similarity ratio, threshold $0.85$) and, for
ranking, the duplicate-penalised \emph{effective accept rate}
$\mathrm{Acc}_{\mathrm{eff}}=\mathrm{Acc}\,(1-\mathrm{DuplicateRate})$.
\textbf{Embedding metrics:} grounding (mean best-match similarity of each question
to the lecture chunks) and diversity ($1$ minus mean nearest-neighbour similarity
between questions). \textbf{Answer metrics:} independent-match rate and
judge-verified-correct rate (Sec.~\ref{sec:bulk-answer}). For the agentic pipeline
only, we additionally report validator--judge agreement ($P$, $R$, $F_1$), treating
the independent judge's accept as ground truth and the self-validator's PASS as the
prediction.

\subsection{Results}
\label{sec:bulk-results}

The tables below report the completed run of $200$ MCQ per pipeline ($400$ total)
for the \emph{agentic} vs.\ \emph{non-agentic (plain)} comparison; the judge is
\texttt{gemma3:12b} and the generator \texttt{gemma3:4b}. \textbf{To illustrate the
full four-pipeline reporting format ahead of the complete run, the not-yet-measured
cells --- the chain-of-thought and program-of-thought rows, the per-dimension judge
scores, diversity, validator precision/recall, and the per-difficulty gold split
--- are shown as illustrative placeholders in \textit{italics}. Italic entries are
example values, not measurements; roman (upright) entries are measured.} The italic
cells are replaced by real figures once the four-pipeline run completes.

\begin{table}[t]
\centering
\small
\caption{Question quality by difficulty and pipeline (judge, $1$--$5$; accept rate
in \%). \emph{Italic} entries are illustrative placeholders (not measured); roman
entries are measured.}
\label{tab:bulk-quality}
\setlength{\tabcolsep}{4pt}
\resizebox{\linewidth}{!}{%
\begin{tabular}{llccccc}
\toprule
Difficulty & Pipeline & Overall & Correctness & Clarity & Diff.\ match & Accept \\
\midrule
Easy   & Agentic         & 4.956 & \textit{4.94} & \textit{4.97} & \textit{4.60} & 100\% \\
       & Non-agentic     & 4.854 & \textit{4.88} & \textit{4.90} & \textit{4.40} & 100\% \\
       & CoT             & \textit{4.90} & \textit{4.92} & \textit{4.93} & \textit{4.55} & \textit{99\%} \\
       & PoT             & \textit{4.91} & \textit{4.95} & \textit{4.90} & \textit{4.58} & \textit{99\%} \\
\midrule
Medium & Agentic         & 4.939 & \textit{4.90} & \textit{4.95} & \textit{4.45} & 99\% \\
       & Non-agentic     & 4.814 & \textit{4.80} & \textit{4.85} & \textit{4.20} & 97\% \\
       & CoT             & \textit{4.88} & \textit{4.90} & \textit{4.90} & \textit{4.40} & \textit{98\%} \\
       & PoT             & \textit{4.90} & \textit{4.93} & \textit{4.88} & \textit{4.50} & \textit{98\%} \\
\midrule
Hard   & Agentic         & 4.884 & \textit{4.88} & \textit{4.90} & \textit{4.20} & 97\% \\
       & Non-agentic     & 4.788 & \textit{4.60} & \textit{4.80} & \textit{4.00} & 91\% \\
       & CoT             & \textit{4.80} & \textit{4.82} & \textit{4.85} & \textit{4.10} & \textit{94\%} \\
       & PoT             & \textit{4.85} & \textit{4.90} & \textit{4.82} & \textit{4.30} & \textit{95\%} \\
\bottomrule
\end{tabular}}
\end{table}

\begin{table}[t]
\centering
\small
\caption{Diversity, grounding, and answer correctness by difficulty and pipeline
(rates in \%; grounding/diversity in $[0,1]$). \emph{Italic} entries are
illustrative placeholders (not measured); roman entries are measured.}
\label{tab:bulk-answer}
\setlength{\tabcolsep}{4pt}
\resizebox{\linewidth}{!}{%
\begin{tabular}{llcccccc}
\toprule
Difficulty & Pipeline & Duplicate & Grounding & Diversity & Eff.\ accept & Indep.\ match & Judge-verified \\
\midrule
Easy   & Agentic         & 35\% & 0.562 & \textit{0.74} & 65\% & 97\% & 100\% \\
       & Non-agentic     & 63\% & 0.642 & \textit{0.55} & 37\% & 95\% & 99\% \\
       & CoT             & \textit{45\%} & \textit{0.60} & \textit{0.66} & \textit{54\%} & \textit{96\%} & \textit{99\%} \\
       & PoT             & \textit{40\%} & \textit{0.61} & \textit{0.69} & \textit{58\%} & \textit{98\%} & \textit{100\%} \\
\midrule
Medium & Agentic         & 38\% & 0.595 & \textit{0.70} & 62\% & 88\% & 99\% \\
       & Non-agentic     & 70\% & 0.614 & \textit{0.48} & 30\% & 84\% & 95\% \\
       & CoT             & \textit{52\%} & \textit{0.60} & \textit{0.60} & \textit{47\%} & \textit{90\%} & \textit{97\%} \\
       & PoT             & \textit{44\%} & \textit{0.62} & \textit{0.66} & \textit{54\%} & \textit{92\%} & \textit{98\%} \\
\midrule
Hard   & Agentic         & 70\% & 0.513 & \textit{0.45} & 30\% & 88\% & 99\% \\
       & Non-agentic     & 38\% & 0.616 & \textit{0.60} & 62\% & 83\% & 81\% \\
       & CoT             & \textit{55\%} & \textit{0.55} & \textit{0.55} & \textit{42\%} & \textit{85\%} & \textit{90\%} \\
       & PoT             & \textit{48\%} & \textit{0.57} & \textit{0.60} & \textit{49\%} & \textit{90\%} & \textit{95\%} \\
\bottomrule
\end{tabular}}
\end{table}

\begin{table}[t]
\centering
\small
\caption{Answer-generator calibration: accuracy on a hand-verified gold set of
$20$ known-answer limits MCQ. \emph{Italic} per-difficulty entries are illustrative
placeholders (not measured); roman overall entries are measured.}
\label{tab:bulk-gold}
\begin{tabular}{llcccc}
\toprule
Model & Role & Overall & Easy & Medium & Hard \\
\midrule
\texttt{gemma3:4b}  & Generator model      & 70\% & \textit{85\%} & \textit{70\%} & \textit{55\%} \\
\texttt{gemma3:12b} & Judge / solver model & 55\% & \textit{70\%} & \textit{55\%} & \textit{40\%} \\
\bottomrule
\end{tabular}
\end{table}

\begin{table}[t]
\centering
\small
\caption{Validator--judge agreement, agentic pipeline. \emph{Italic}
precision/recall are illustrative placeholders (not measured); roman $F_1$ is
measured.}
\label{tab:bulk-f1}
\begin{tabular}{lccc}
\toprule
Difficulty & Precision & Recall & $F_1$ \\
\midrule
Easy   & \textit{1.00} & \textit{0.65} & 0.786 \\
Medium & \textit{0.97} & \textit{0.64} & 0.774 \\
Hard   & \textit{1.00} & \textit{0.73} & 0.842 \\
\bottomrule
\end{tabular}
\end{table}

\paragraph{Findings.} Three observations follow from
Tables~\ref{tab:bulk-quality}--\ref{tab:bulk-gold}.
\begin{enumerate}
  \item \textbf{The agentic pipeline is more \emph{answer-correct}, most clearly
        on hard questions.} Judge-verified correctness is $99\%$ (agentic) vs
        $81\%$ (non-agentic) at the hard level, and independent re-solve accuracy
        is higher at every difficulty ($97/88/88\%$ vs $95/84/83\%$). Mean judge
        score and accept rate also favour the agentic pipeline throughout.
  \item \textbf{Diversity is the discriminator, and it reverses at hard.} At easy
        and medium the agentic pipeline is far less repetitive ($35\%/38\%$
        duplicates vs $63\%/70\%$), but at hard it \emph{regresses} to $70\%$
        duplicates against the baseline's $38\%$. Because the effective accept rate
        discounts repetition, this single metric flips the hard-level ranking:
        $30\%$ (agentic) vs $62\%$ (non-agentic). The agentic pipeline concentrates
        hard questions on a narrow band of the material and repeats itself there.
  \item \textbf{The answer-checker is only moderately reliable on this domain, so
        answer-accuracy is agreement, not verified truth.} On the gold set the
        independent solver (\texttt{gemma3:12b}) scores $55\%$ and the generator
        model (\texttt{gemma3:4b}) $70\%$. High generator--solver agreement
        ($88$--$97\%$) therefore partly reflects shared model behaviour rather than
        confirmed correctness; the independent-match rate should be read alongside
        this calibration. Strengthening the solver (e.g.\ program-of-thought
        derivation of each limit) is the direct lever for raising it.
\end{enumerate}

\subsection{Implementation status}
\label{sec:bulk-status}

The harness described above is implemented and runnable: the four pipelines, the
independent Answer Generator, the gold cross-check, per-difficulty and overall
aggregation, per-question CSV/JSON export, and charts. The batch runs as a
resumable background job with a live progress feed, so a disconnect resumes from
the last completed question rather than restarting. The agentic-vs-non-agentic run
of Sec.~\ref{sec:bulk-results} ($400$ MCQ) is complete; the chain-of-thought and
program-of-thought columns are populated when the four-pipeline run is executed on
the same source chapter.

\paragraph{Limitations.} The judge and Answer Generator are LLMs rather than human
raters, so absolute scores are a relative signal; the generator and judge share the
Gemma-3 family (mitigated, not eliminated, by using a larger independent judge for
scoring and solving); and figures are from a single source chapter, so per-cell
values should be read with the corresponding $N$ in mind.
